% Compile with:  lualatex ctr_structure_factor_physics.tex   (twice, for refs)
\documentclass[11pt,a4paper]{article}

\usepackage[T1]{fontenc}
\usepackage{lmodern}
\usepackage[margin=27mm]{geometry}
\usepackage{amsmath}
\usepackage{amssymb}
\usepackage{booktabs}
\usepackage{array}
\usepackage{enumitem}
\usepackage{xcolor}
\usepackage[colorlinks=true,linkcolor=blue!45!black,urlcolor=blue!45!black,
            citecolor=blue!45!black]{hyperref}

\newcommand{\Fhkl}{\left|F_{hkl}\right|^{2}}
\newcommand{\dgam}{\Delta\gamma}
\newcommand{\dOm}{\Delta\Omega}
\newcommand{\Ored}{\widetilde{\Omega}}
\newcommand{\Crod}{C_{\mathrm{rod}}}
\newcommand{\Cdet}{C_{\mathrm{det}}}
\newcommand{\Carea}{C_{\mathrm{area}}}
\newcommand{\Cbeam}{C_{\mathrm{beam}}}
\newcommand{\Cillum}{C_{\mathrm{illum}}}
\newcommand{\bin}{\beta_{\mathrm{in}}}
\newcommand{\bout}{\beta_{\mathrm{out}}}
\newcommand{\re}{r_{e}}
\newcommand{\Au}{A_{u}}
% Long dotted module paths are unbreakable in \texttt; let them break after
% a dot, and give TeX a little slack rather than an overfull line.
\newcommand{\code}[1]{\texttt{\small #1}}
\newcommand{\dt}{.\allowbreak}
\setlength{\emergencystretch}{3em}

\newcommand{\unit}[1]{\;\mathrm{#1}}
\newcommand{\Ang}{\text{\AA}}

\title{Structure factors from area-detector surface diffraction data\\
       \large The normalizations and integration intervals used by orGUI}
\author{orGUI --- \code{orgui\dt datautils\dt xrayutils\dt corrections}}
\date{\today}

\begin{document}
\maketitle

\begin{abstract}
\noindent
This document states, in one place, how orGUI turns detector counts into
$\Fhkl$: every normalization that is divided out, every integration interval
that is summed or integrated over, and the unit each quantity carries at the
boundary where it changes. It covers the two direct-space measurement modes
--- rocking scans and stationary area-detector scans --- and the conversion to
and from specular reflectivity. The reciprocal-space reconstruction is a third
route with a different set of factors and is only contrasted here, not
derived.

The physics is that of E.~Vlieg, \emph{J.~Appl.~Cryst.} \textbf{30} (1997) 532,
in the two-dimensional-detector form of J.~Drnec \emph{et al.},
\emph{J.~Appl.~Cryst.} \textbf{47} (2014) 365. What this document adds is the
mapping from those expressions onto what the code actually sums, and the
resolution of two couplings that the published expressions do not address
because they do not arise with a point detector: the interaction between the
detector solid-angle correction and the out-of-plane acceptance, and the
choice of integration variable and its interval.
\end{abstract}

\tableofcontents

\section{Conventions}

\subsection{Geometry and angles}

orGUI works in the \textbf{z-axis geometry}, with angles as defined in
\code{orgui\dt datautils\dt xrayutils\dt HKLVlieg}. It is essential to
separate the \emph{motor positions} of the diffractometer from the
\emph{scattering angles} of an individual detector pixel, because only the
latter enter any correction factor:

\begin{center}
\begin{tabular}{@{}llll@{}}
\toprule
Symbol & Meaning & Kind & Varies with \\
\midrule
$\alpha$ & incidence angle on the surface & circle (\code{mu})
  & frame \\
$\omega$ & sample rotation about the surface normal & circle (\code{th})
  & frame \\
$\gamma_{\mathrm{arm}}$, $\delta_{\mathrm{arm}}$
  & detector arm position & circles & frame \\
\addlinespace
$\gamma_p$, $\delta_p$ & laboratory-frame angles of a pixel & derived
  & pixel, arm \\
$\gamma$, $\delta$ & \textbf{surface-frame} angles of a pixel & derived
  & pixel, arm, $\alpha$ \\
\bottomrule
\end{tabular}
\end{center}

\noindent
$\gamma$ and $\delta$ are \textbf{not} diffractometer circles and are
\textbf{not} the arm angles. They are per-pixel quantities, computed from
where a pixel sits \emph{after} the detector has been rotated to the arm
position of that frame. Within one frame they vary across the detector face;
between frames they change both because the arm moved and because $\alpha$
changed. In this geometry $\gamma$ is the exit angle of the diffracted beam
from the surface, so $\bout = \gamma$ and $\bin = \alpha$.

All angles are in \textbf{radian} everywhere inside
\code{orgui\dt datautils\dt xrayutils\dt corrections}; degrees appear only at
the user interface and in raw motor positions read from a scan.

\subsection{From a pixel to its scattering angles}
\label{sec:pixelangles}

The conversion runs in three steps, all in
\code{DetectorCalibration\dt Detector2D\_SXRD}.

\paragraph{1. The arm position becomes a detector geometry.} The arm angles
are given as the true laboratory-frame angles
$(\gamma_{\mathrm{arm}}, \delta_{\mathrm{arm}})$ the arm reference direction
points at. The rotation from the \emph{calibrated} arm position to the
requested one,
\begin{equation}
  R_{\mathrm{rel}} = R_{\mathrm{arm}}(\gamma_{\mathrm{arm}},
                                      \delta_{\mathrm{arm}})\,
                     R_{\mathrm{arm}}^{\mathrm{ref}\,\top} ,
\end{equation}
is folded into the pyFAI \emph{home} geometry
$[\,d, \mathrm{poni}_1, \mathrm{poni}_2, \mathrm{rot}_1, \mathrm{rot}_2,
\mathrm{rot}_3\,]$ by \code{paramAtArm}. A moving arm is therefore represented
as a \emph{modified pyFAI geometry}, not as an angular offset added after the
fact --- which is what makes the parallax and obliquity of the following step
come out right at every arm position.

Two conventions here are easy to get wrong. Passing \code{None} for both arm
angles means \emph{at the calibrated arm position}, which is not the same as
passing $0.0$ unless the calibration was taken with the motors reading zero.
And exactly one of the two is rejected rather than defaulted, because the
calibration reference is a rotation and not a pair of independent offsets.

\paragraph{2. Pixel to laboratory frame.} On that geometry, the pixel position
gives the scattering angle $2\theta$ and the azimuth $\chi$ --- with parallax
applied, and $\chi$ offset by the azimuthal reference --- and from those the
laboratory-frame angles $(\gamma_p, \delta_p)$ of the pixel
(\code{primBeamPoints}).

\paragraph{3. Laboratory to surface frame.} Tilting into the surface frame of
a sample at incidence angle $\alpha$ (\code{surfaceAnglesPoint}):
\begin{align}
  \gamma &= \arcsin\bigl(\cos\alpha \, \sin\gamma_p
            - \sin\alpha \, \cos\delta_p \, \cos\gamma_p \bigr) ,
  \label{eq:gamma}\\[2pt]
  \delta &= \arcsin\!\left(
            \frac{\sin\delta_p \, \cos\gamma_p}{\cos\gamma} \right) .
  \label{eq:delta}
\end{align}
At $\alpha = 0$ these collapse to $\gamma = \gamma_p$ and
$\delta = \delta_p$, as they must: at zero incidence the surface frame is the
laboratory frame.

\paragraph{Which angles the corrections use.} Every factor in this document
that names $\gamma$ or $\delta$ means the surface-frame angles of
Eqs.~\eqref{eq:gamma} and~\eqref{eq:delta}, evaluated \textbf{at the
reflection position} --- the centre of the region of interest --- for the arm
position and $\alpha$ of that frame. The integration paths read them from the
stored per-reflection angles, not from arm motors. Two places in this document
depend on the distinction directly:

\begin{itemize}[leftmargin=1.4em]
\item $\dgam$ (Section~\ref{sec:acceptance}) is a \emph{difference} of
  Eq.~\eqref{eq:gamma} between two pixel rows, which is only meaningful
  because $\gamma$ varies across the detector face within one frame.
\item The polarization factor (Section~\ref{sec:notmodelled}) is built on the
  \emph{home} geometry rather than at the frame's arm position, which is why
  it needs a per-frame correction for a scan that drives the arm --- and why
  the reconstruction, which cannot use that correction, is still wrong there.
\end{itemize}

\subsection{Units at the boundaries}

Units are deliberately mixed, following the repository convention that
crystallographic quantities are in \AA ngstrom and instrument quantities in
meter. Getting this wrong is a factor of $10^{20}$ in the scale factor, so it
is stated at every function boundary and pinned by
\code{test\_scale\_factor\_carries\_its\_documented\_units}.

\begin{center}
\begin{tabular}{@{}lll@{}}
\toprule
Quantity & Symbol & Unit \\
\midrule
Wavelength                    & $\lambda$ & \AA ngstrom \\
Surface unit-cell area        & $\Au$     & \AA ngstrom squared \\
Classical electron radius     & $\re$     & meter ($2.8179403262\times10^{-15}$) \\
Active illuminated area       & $A$       & square meter \\
Incident flux density         & $\Phi_0$  & photons per second and square meter \\
Counting time                 & $T$       & second \\
All angles                    &           & radian \\
Structure factor squared      & $\Fhkl$   & electron units squared \\
\bottomrule
\end{tabular}
\end{center}

\noindent
$\Phi_0$ is a flux \emph{density}: $\Phi_0 A_0$ is the total flux on the
sample, with $A_0$ the beam cross-section.

\section{What one detector frame contains}
\label{sec:frame}

A pixel $i$ of one frame of counting time $T$ accumulates
\begin{equation}
  N_i = \Phi_0\, T \left(\frac{d\sigma}{d\Omega}\right)_{\!i} \dOm_i ,
  \label{eq:pixel}
\end{equation}
where $\dOm_i$ is the solid angle that pixel subtends at the sample. For a
flat detector at distance $d$ with square pixels of size $p$, seen at an
incidence angle $\theta$ on the detector face,
\begin{equation}
  \dOm = \frac{p^{2}}{d^{2}}\cos^{3}\theta ,
  \qquad
  \delta\gamma = \frac{p}{d}\cos^{2}\theta ,
  \qquad
  \delta\psi = \frac{p}{d}\cos\theta ,
  \label{eq:pixelgeom}
\end{equation}
with $\delta\gamma$ the angular height of the pixel in the direction it is
tilted away from the normal and $\delta\psi$ the angular width perpendicular
to it; their product is $\dOm$ as it must be. pyFAI's
\code{solidAngleArray} returns the \emph{normalized} solid angle
\begin{equation}
  \Ored = \dOm \Big/ \frac{p^{2}}{d^{2}} = \cos^{3}\theta ,
\end{equation}
which equals $1$ at the point of normal incidence and which orGUI verifies
against $\cos^{3}\theta$ to $10^{-9}$ in
\code{test\_corrections\_acceptance.py}.

Equation~\eqref{eq:pixel} is the origin of the central result of
Section~\ref{sec:solidangle}: a \emph{sum} of $N_i$ over a region already
carries each pixel's own $\dOm_i$, and is therefore already an angular
integral. Only a quantity that is meant to be a differential cross-section at
a point needs $\dOm_i$ divided out.

\section{Integration intervals}
\label{sec:intervals}

This section fixes exactly which counts enter the integrated intensity, for
each mode. Everything here is a statement about \emph{domains}; the divisors
follow in Section~\ref{sec:normalizations}.

\subsection{The region of interest, and its background}

Both direct-space modes reduce a frame to a single number by summing a
rectangular signal region and subtracting a background estimated from
separate background regions.

\paragraph{Stationary mode} (\code{orGUI.integrateROI}). With $S$ the signal
region, $B$ the union of the background regions, $n_S$ and $n_B$ the numbers
of \emph{valid} (unmasked, present) pixels in each, and $n_S^{\mathrm{nom}}$
the nominal region size in pixels,
\begin{equation}
  N \;=\; \Biggl(\underbrace{\sum_{i\in S} N_i}_{\text{signal sum}}
        \;-\; \frac{n_S}{n_B}\underbrace{\sum_{i\in B} N_i}_{\text{background sum}}
    \Biggr)\,
    \frac{n_S^{\mathrm{nom}}}{n_S} .
  \label{eq:stationary_roi}
\end{equation}
The factor $n_S/n_B$ puts the background on a per-pixel basis and scales it to
the signal region. The trailing $n_S^{\mathrm{nom}}/n_S$ rescales from the
pixels that were actually valid to the nominal region area, so that masked
pixels and detector gaps do not reduce the integrated intensity.

That last rescaling is the reason the nominal region area must \emph{not}
appear again in any correction factor: the ROI-mean bookkeeping of
\code{corrections.roi.roi\_mean\_correction} divides the summed correction by
the valid-pixel count and never by the nominal area. Multiplying by the area a
second time would scale every intensity by the size of its own region, and
because orGUI resizes regions with detector position
(\code{ROIutils.calc\_corrections}), two measurements of one rod taken on
different parts of the detector would be scaled apart.

\paragraph{Rocking mode} (\code{peak1Dintegr}). The same per-frame reduction is
applied first, giving a curve $N(\omega_k)$ over the frames of the scan, and
that curve is then integrated over the rocking angle.

\subsection{The rocking integral}
\label{sec:rockinginterval}

For each output point the integration interval is a closed interval
$[\omega_{\mathrm{from}}, \omega_{\mathrm{to}}]$, snapped to the nearest
sampled motor positions and ordered so that the result is positive:
\begin{equation}
  \mathcal{I} \;=\; \int_{\omega_{\mathrm{from}}}^{\omega_{\mathrm{to}}}
                     N(\omega)\, d\omega
    \;\longrightarrow\;
    \sum_{k=k_{\mathrm{from}}}^{k_{\mathrm{to}}-1}
      \tfrac{1}{2}\bigl(N_k + N_{k+1}\bigr)\,(\omega_{k+1}-\omega_k) ,
\end{equation}
a trapezoidal rule over a slice that \textbf{includes both endpoints}
(\code{slice(idx\_from, idx\_to + 1)}), so that the integrated samples, the
error weights and the reported interval width all describe the same domain.
The interval width recorded alongside the result is
$\Delta\omega = |\omega_{k_{\mathrm{to}}} - \omega_{k_{\mathrm{from}}}|$.

Background regions are integrated over their own intervals and then scaled to
the signal interval by the ratio of interval widths,
$\Delta\omega_S/\sum_{B}\Delta\omega_B$, before being subtracted. The result
is deliberately \emph{not} divided by $\Delta\omega_S$: it is an integral, not
a mean, because that is what Vlieg's expression requires.

\paragraph{The integration variable must be in radian.} The published
expressions integrate $d\omega$ in radian. The motor axis is in degrees, so
the integral as computed is larger by $180/\pi = 57.3$. This is a constant, so
it is invisible within one data set and only appears when a rocking scan is
compared with a stationary scan, or with an absolute scale.

\subsection{The out-of-plane acceptance $\dgam$}
\label{sec:acceptance}

A rocking scan does not measure a whole rod. It intercepts a slice whose
length is (Vlieg eq.~20)
\begin{equation}
  \Delta l = \Crod \, \frac{V_u}{\lambda \Au}\, \dgam ,
\end{equation}
so its integrated intensity is proportional to $\dgam$, the out-of-plane
angular acceptance of the region of interest. A stationary measurement
intercepts the whole rod cross-section and carries no such factor.

With a point detector $\dgam$ was a slit setting: fixed for an experiment and
absorbed into an overall scale factor. On an area detector it is a property of
the region, and orGUI sizes regions per detector position, so it varies
\emph{along a scan} and changes the shape of a rod rather than only its scale.

\code{corrections.acceptance} evaluates it as the exit-angle span between the
top and bottom \textbf{edges} of the region, at its centre column:
\begin{equation}
  \dgam = \Bigl|\,
    \gamma\bigl(r + \tfrac{n}{2},\, c\bigr) -
    \gamma\bigl(r - \tfrac{n}{2},\, c\bigr)\Bigr| ,
  \label{eq:dgamma}
\end{equation}
for a region of $n$ rows centred on pixel $(r,c)$, with $\gamma(\cdot)$ the
surface-frame exit angle of Eq.~\eqref{eq:gamma} evaluated at that pixel and
at the frame's arm position and $\alpha$. Two choices in
Eq.~\eqref{eq:dgamma} are deliberate:

\begin{itemize}[leftmargin=1.4em]
\item \textbf{Edges, not pixel centres.} A region of $n$ rows accepts $n$
  pixels' worth of angle, not $n-1$. Measuring between the centres of the
  first and last row understates every acceptance by one pixel: $1\,\%$ on a
  100-row region, $10\,\%$ on a 10-row one.
\item \textbf{The centre column.} The region is centred on the rod, and it is
  the range of $\gamma$ over which \emph{the rod} crosses the aperture that
  sets the intercepted length. The corner-to-corner span over the whole
  rectangle is computed separately by \code{gamma\_range}; the two agree when
  lines of constant $\gamma$ run along the detector rows and separate on a
  detector rolled about the beam, so their ratio is a cheap diagnostic for
  whether the estimate can be trusted.
\end{itemize}

\paragraph{The detector arm barely matters --- for the span.} Driving the
$\gamma$ arm shifts $\gamma$ at the region centre by the full arm angle, so
the arm position matters a great deal for \emph{where on the rod} the region
sits. It leaves the \emph{span} of Eq.~\eqref{eq:dgamma} unchanged to nine
digits, because that rotation is about the very axis $\gamma$ is measured
around; a $40^{\circ}$ $\delta$-arm rotation moves it by one part in $10^{5}$.
So how much rod is accepted is nearly independent of the arm even though which
part is accepted is not, and an acceptance computed without arm bookkeeping
remains usable. Pass the arm angles anyway where they are known --- they cost
nothing. This is the opposite of the polarization factor
(Section~\ref{sec:notmodelled}), where ignoring the arm is a $10\,\%$ error at
a scattering angle of $18^{\circ}$.

\section{Normalizations}
\label{sec:normalizations}

\subsection{Counting time and monitor}

Both published expressions carry $\Phi_0 T$, so neither an integrated rocking
curve nor a stationary frame means anything until it is divided by its
counting time and by whatever counter the incident flux is tracked with:
\begin{equation}
  I = \frac{\mathcal{N}}{T \prod_j M_j} ,
\end{equation}
with $\mathcal{N}$ the background-subtracted integrated counts of
Section~\ref{sec:intervals} and $M_j$ the monitor counters
(\code{corrections.normalization.normalization\_divisor}).

For a rocking scan this is the \textbf{per-frame} counting time, not the sum
over the scan. The rocking angle is the integration variable; the time is not.
Using the summed time would divide by the number of frames a second time.

Beyond mode equivalence, this matters \emph{within} one data set: rocking
scans taken at different $l$ with different counting times, or with a drifting
ring current, are otherwise placed on different scales with nothing in the
saved data recording it.

\subsection{Polarization}

The polarization factor $P$ is a per-pixel property of the scattering angle.
It is divided out as the mean of $1/P$ over the region
(\code{roi\_mean\_correction}) before the reduction; the master equation of
Section~\ref{sec:master} therefore takes an intensity from which $P$ has
already been removed.

\subsection{Solid angle: why a summed region takes none}
\label{sec:solidangle}

This is the one normalization whose \emph{absence} needs an argument, because
the array is available and applying it looks harmless.

\paragraph{The statement.} A region-summed intensity gets no solid-angle
correction, in either direct-space mode. The array belongs only to the
per-pixel reciprocal-space reconstruction, which forms a differential
cross-section rather than a sum over an aperture.

\paragraph{Why.} From Eq.~\eqref{eq:pixel}, the raw sum over a region is
\begin{equation}
  \sum_{i\in S} N_i
  = \Phi_0 T \sum_{i \in S} \left(\frac{d\sigma}{d\Omega}\right)_{\!i} \dOm_i
  \;\simeq\; \Phi_0 T \iint_{S} \frac{d\sigma}{d\Omega}\, d\gamma\, d\psi ,
\end{equation}
already the complete angular integral over the aperture. Every photon is
counted once, weighted by the solid angle of the pixel that caught it. There
is nothing to correct. Three remarks make this robust:

\begin{enumerate}[leftmargin=1.6em]
\item If the reflection were \emph{not} fully contained --- post-sample slits
  narrower than the peak --- the missing part is Vlieg's $\Cdet$, an
  aperture-against-profile model. A per-pixel obliquity weight does not repair
  a truncated tail, so the conclusion does not depend on the open-slit
  assumption holding.
\item A rocking sum is a complete integral too, over a $\dgam$ slice of the
  rod. The slice is what $\dgam$ accounts for; there is still no per-pixel
  weighting to undo.
\item Where a differential cross-section \emph{is} the goal --- binning
  individual pixels into reciprocal-space voxels --- the correction is
  required, and the reconstruction path applies it under its own switch.
\end{enumerate}

\paragraph{What goes wrong if it is applied anyway.} It does not merely add a
constant. For a rocking scan the $\omega$-integrated counts in pixel row $j$
go as that row's angular height $\delta\gamma_j$, so
\begin{equation}
  \text{raw sum} \;\longleftrightarrow\; \sum_j \delta\gamma_j = \dgam ,
  \qquad
  \text{corrected sum} \;\longleftrightarrow\;
      \sum_j \frac{\delta\gamma_j}{\Ored_j} \simeq \frac{\dgam}{\langle\Ored\rangle} .
  \label{eq:partners}
\end{equation}
The solid-angle factor cancels against the acceptance in the rocking mode ---
which is exactly why leaving it in \emph{looks} harmless there. A stationary
sum has no $\dgam$ for it to cancel against, so the factor survives into the
saved number and the two modes disagree by $\langle 1/\Ored\rangle$: $0.7\,\%$
for a detector at $1$~m and $7\,\%$ at $0.3$~m.

Table~\ref{tab:f6} measures Eq.~\eqref{eq:partners} on a calibrated
$172\,\mu$m detector with a 60-row region. The last two columns
agreeing to $5\times10^{-6}$ is the cancellation.

\begin{table}[htbp]
\centering
\begin{tabular}{@{}llrrrr@{}}
\toprule
$d$ & $\theta$ & $\dgam/\dgam_{\mathrm{nom}}$ & exact$/\dgam_{\mathrm{nom}}$
    & exact$/\dgam$ & $1/\langle\Ored\rangle$ \\
\midrule
$1.0$~m & $2.7^{\circ}$ & 0.99774 & 1.00112 & 1.003383 & 1.003383 \\
$0.3$~m & $5.6^{\circ}$ & 0.99039 & 1.00476 & 1.014512 & 1.014509 \\
$0.3$~m & $9.0^{\circ}$ & 0.97549 & 1.01236 & 1.037795 & 1.037788 \\
\bottomrule
\end{tabular}
\caption{The two partners of Eq.~\eqref{eq:partners}, against the nominal
$\dgam_{\mathrm{nom}} = n\,p/d$. Here the region is displaced from the point
of normal incidence along its rows, so by Eq.~\eqref{eq:pixelgeom}
$\dgam = \dgam_{\mathrm{nom}}\cos^{2}\theta$ against $\Ored = \cos^{3}\theta$
and the corrected-sum partner is $\dgam_{\mathrm{nom}}/\cos\theta$. That
decomposition depends on the orientation of the region relative to the normal;
the ratio in the last two columns does not, since it follows from
Eq.~\eqref{eq:partners} alone. Pinned by
\code{test\_the\_solid\_angle\_correction\_and\_the\_acceptance\_carry\_one\_obliquity}.}
\label{tab:f6}
\end{table}

\paragraph{Two traps this leaves behind.}
\begin{itemize}[leftmargin=1.4em]
\item \textbf{$\dgam$ does not become unnecessary.} It is a \emph{size} factor
  --- how many rows of rod the region accepted --- while the solid angle is an
  \emph{obliquity} factor. They overlap only in the obliquity, a fraction of a
  percent.
\item \textbf{A constant nominal $n\,p/d$ is not a valid substitute for
  Eq.~\eqref{eq:dgamma}.} It is wrong by $1/\cos\theta$: $0.1\,\%$ at $1$~m,
  $1.2\,\%$ at $\theta = 9^{\circ}$. The single-pixel acceptance times the row
  count is the same angle-derived quantity as Eq.~\eqref{eq:dgamma}, not a
  cheaper alternative to it.
\end{itemize}

\subsection{The active area and the beam footprint}
\label{sec:area}

$A$ is the surface area that actually contributes counts. Vlieg writes it as
$A = A_0\,\Carea\,\Cbeam$ (his eq.~41). Which of two limits applies is an
experimental question, and orGUI builds both:

\paragraph{Beam-limited (open slits, orGUI's usual case).}
\begin{equation}
  A(\alpha) = w\, L\, \Cillum(\alpha, L) ,
  \label{eq:beamarea}
\end{equation}
with $w$ the horizontal beam width, $L$ the sample length along the beam, and
$\Cillum$ the profile integrated over the projected sample and normalized to a
uniform beam,
\begin{equation}
  \Cillum(\alpha,L)
    = \frac{1}{p_{\max}\, h_{\mathrm{proj}}}
      \int_{-h_{\mathrm{proj}}/2}^{+h_{\mathrm{proj}}/2} p(z)\, dz ,
  \qquad h_{\mathrm{proj}} = L\sin\alpha .
\end{equation}
Note that $\Cillum$ is a dimensionless \emph{fraction}, not an area; the
horizontal direction is the plain beam width, since it is not projected. For a
top-hat profile of full height $h$, Eq.~\eqref{eq:beamarea} reduces exactly to
the familiar closed form
\begin{equation}
  A = w \min\!\left(L,\; \frac{h}{\sin\alpha}\right) ,
\end{equation}
verified to $3.3\times10^{-16}$. It does \emph{not} reduce to it for any other
profile --- a Gaussian of the same width departs by up to $19\,\%$ --- which
is why the closed form is not used.

\paragraph{Slit-limited.} When post-sample slits cut the footprint,
$\Carea = 1/(\sin\delta\,\cos(\alpha - \bin))$ applies instead (Vlieg eqs.~37
and~38). Because $\delta$ varies along a rod, this is a \emph{shape} effect,
not only a scale: on a simulated Pt(111) $(1,0,l)$ rod, $\delta$ moves from
$16.77^{\circ}$ to $16.95^{\circ}$ between $l = 0.4$ and $l = 3.0$, a
$1.1\,\%$ error, and much larger between rods at different in-plane momentum
transfer.

\paragraph{It cannot explain a mode disagreement.} $A$ enters the rocking and
the stationary expressions identically and cancels from their ratio. It sets
the absolute scale and nothing else.

\section{The master equation}
\label{sec:master}

With $I$ the normalized integrated intensity of
Section~\ref{sec:normalizations} --- counts per second per monitor unit, with
$P$ already divided out ---
\begin{equation}
  \boxed{\;
  I \;=\; \Phi_0\, \frac{\re^{2} A \lambda^{2}}{\Au^{2}}\;
          \Fhkl \; \eta \; \Cdet \; }
  \label{eq:master}
\end{equation}
and the mode enters only through the angular factor $\eta$:

\begin{center}
\begin{tabular}{@{}lll@{}}
\toprule
Mode & $\eta$ & Lorentz factor \\
\midrule
Rocking scan (\code{th}/$\omega$)      & $L_\varphi\,\Crod\,\dgam$
  & $L_\varphi = 1/(\sin\delta\,\cos\alpha\,\cos\gamma)$ \\
Reflectivity rocking (\code{mu})       & $L_r\,\Crod\,\dgam$
  & $L_r = 1/\sin 2\alpha$ \\
Stationary area detector               & $L_s$
  & $L_s = 1/\sin\gamma$ \\
Specular reflectivity                  & $L_s$ at $\gamma = \alpha$
  & $L_s = 1/\sin\alpha$ \\
\bottomrule
\end{tabular}
\end{center}

\noindent
with the rod interception $\Crod = \cos\gamma$. The three Lorentz factors are
\textbf{alternatives selected by how the intensity was measured}, not factors
to be combined. Two structural points:

\begin{itemize}[leftmargin=1.4em]
\item The rocking modes carry $\Crod$ and $\dgam$; the stationary mode carries
  neither. A stationary measurement integrates across the whole rod
  cross-section, so there is no interception factor and no acceptance ---
  \code{angular\_factor} \emph{rejects} an acceptance passed for a stationary
  mode, rather than ignoring it.
\item The area correction $1/\sin\delta$ listed in the ANA/ROD z-axis table is
  \emph{not} applied; the numerically evaluated footprint of
  Section~\ref{sec:area} is used instead, which is the row the manual marks as
  calculated numerically.
\end{itemize}

\subsection{Dimensional check}

The radian in $\dgam$ is what makes the two modes dimensionally distinct, and
it is worth checking that it lands correctly. Writing $[\cdot]$ for the unit
of a quantity and taking a unit monitor,
\begin{align}
  \text{prefactor:}&\quad
    \frac{1}{\mathrm{s}\,\mathrm{m}^{2}} \cdot \mathrm{m}^{2}
    \cdot \mathrm{m}^{2} \cdot \frac{\mathrm{m}^{2}}{\mathrm{m}^{4}}
    = \frac{1}{\mathrm{s}} , \\
  \text{stationary:}&\quad
    [I] = \frac{\text{counts}}{\mathrm{s}} ,
    \qquad \eta \text{ dimensionless} \quad\checkmark \\
  \text{rocking:}&\quad
    [I] = \frac{\text{counts}\cdot\mathrm{rad}}{\mathrm{s}} ,
    \qquad [\eta] = \mathrm{rad} \quad\checkmark
\end{align}
The rocking intensity carries one radian more than the stationary one, and
$\dgam$ is what absorbs it. Integrating in degrees and omitting $\dgam$ are
therefore not two independent mistakes but a single unbalanced unit, which is
why they were measured together as one ratio.

\subsection{Inverting to a structure factor}

\begin{equation}
  \Fhkl = \frac{I}{S\,\eta\,\Cdet} ,
  \qquad
  S = \Phi_0\, A\, \frac{\re^{2}\lambda^{2}}{\Au^{2}} ,
\end{equation}
with $S$ the mode-independent prefactor, in $1/\mathrm{s}$
(\code{measurement.scale\_factor}). Leaving $\Phi_0 = A = 1$ gives the
\textbf{relative} scale, which is all that is needed to place different scan
modes of one experiment on a common scale. Supplying the measured flux density
and illuminated area gives the \textbf{absolute} scale, and $\Fhkl$ then comes
out in electron units squared --- the same scale on which
\code{CTRcalc.SXRDCrystal} calculates.

\section{Reflectivity}

Specular reflectivity is the stationary case at $\gamma = \alpha$ with a
beam-limited active area $A_r = A_0/\sin\alpha$. Vlieg's eq.~63, generalized to
a non-specular exit angle, is
\begin{equation}
  R = \frac{\re^{2}\lambda^{2} P_r}{\Au^{2}\,\sin\alpha\,\sin\bout}\, \Fhkl ,
  \label{eq:refl}
\end{equation}
the fraction of the incident flux scattered into the rod. For the specular rod
$\bout = \alpha$ and this is the familiar $1/\sin^{2}\alpha$: one power from
the footprint, one from the stationary Lorentz factor.

Nothing beyond an absolutely scaled $\Fhkl$ is required, so \textbf{an absolute
reflectivity comes for free} once the scale factor is absolute, and a
reflectivity curve can be refined together with the truncation rods.
Equation~\eqref{eq:refl} inverts to put a measured reflectivity onto the
$\Fhkl$ scale.

Three caveats:
\begin{enumerate}[leftmargin=1.6em]
\item It is the \textbf{kinematic} result. It fails near a bulk Bragg peak and
  below the critical angle, where refraction and multiple scattering dominate
  --- which is where much of the interesting part of a reflectivity curve
  lies. Compare against a DWBA treatment there.
\item A reflectivity scan \textbf{drives the detector arm}, which is where the
  polarization issue of Section~\ref{sec:notmodelled} bit hardest. The two
  direct-space integration paths correct for it; a reflectivity reduced from a
  reconstructed map does not.
\item \textbf{Off-specular}, $R$ is well defined for a truncation rod
  integrated across its cross-section, but not for diffuse scattering, where
  only a differential cross-section is meaningful.
\end{enumerate}

Note also that the ANA ``reflectivity rocking scan'' Lorentz factor
$1/\sin 2\alpha$ describes rocking \emph{through} the specular ridge. That is a
different measurement from a stationary specular scan; both are legitimate and
they are not interchangeable.

\section{What is deliberately not modelled}
\label{sec:notmodelled}

\begin{description}[leftmargin=0pt,style=unboxed]
\item[$\Cdet$, the in-plane detector acceptance.] Vlieg's sections 2.4--2.5.
  It is close to $1$ for a region wide enough to contain the whole in-plane
  peak profile, which is what Eq.~\eqref{eq:master} assumes by taking
  $\Cdet = 1$. It is \emph{not} the same for the two modes when the rod is
  broad: Drnec's Fig.~18 shows direct stationary integration underestimating
  $|F|$ at low $l$ by up to a factor of $2$ for exactly this reason. This is
  the one remaining mechanism that can make the two modes disagree after
  everything in this document is applied, and it is a modelling problem rather
  than a normalization one. It cannot be validated on simulated data --- it
  needs the overlap region of a rocking and a stationary scan on the same rod.

\item[The polarization at a moving detector arm, in the reconstruction.] The
  per-pixel polarization array is built once, outside the frame loop, from the
  \emph{home} geometry --- step 1 of Section~\ref{sec:pixelangles} is skipped,
  so it is the array at the calibrated arm position rather than at the
  frame's. For a fixed arm that is correct: each pixel already carries its own
  scattering angle, and the apparent $\alpha$ dependence of the z-axis
  expression is only a change of frame. For a scan that drives the arm it
  understates the correction: $3\,\%$ at a scattering angle of $10^{\circ}$,
  $10\,\%$ at $18^{\circ}$ and $33\,\%$ at $30^{\circ}$. This is the opposite
  of the acceptance (Section~\ref{sec:acceptance}), whose \emph{span} survives
  a moving arm unchanged, and of the solid angle, which is invariant under it
  exactly --- an arm rotation is a rigid rotation about the sample, so every
  pixel keeps its distance and its obliquity to its own line of sight.

  The two direct-space integration paths \textbf{correct} for this: they
  reduce the polarization to a region mean anyway, so a per-frame ratio of two
  region means, $\langle 1/P_{\mathrm{arm}}\rangle /
  \langle 1/P_{\mathrm{home}}\rangle$, moves it onto the frame's arm position.
  That ratio is exactly $1$ at the calibrated position, so a fixed-arm scan is
  untouched. What remains uncorrected is the \textbf{reciprocal-space
  reconstruction}, which applies the polarization per pixel rather than as a
  region mean and would need the array itself rebuilt per frame.

\item[Uncertainty of $\dgam$.] Errors follow the same divisors as the
  intensities, but a $\dgam$ estimated from the calibrated geometry has an
  uncertainty of its own that is not currently propagated.

\item[The reciprocal-space route.] Voxel binning absorbs the Lorentz factor
  (Drnec section 4), so a reconstructed map needs $\Carea$, $\Cbeam$, $P$ and
  a $\Delta l$, but no Lorentz factor --- a different set of factors from
  either direct-space mode. Placing it on the same absolute scale is separate
  work.
\end{description}

\section{Summary}

\begin{center}
\renewcommand{\arraystretch}{1.18}
\begin{tabular}{@{}lllcc@{}}
\toprule
Factor & Symbol & Where & Rocking & Stationary \\
\midrule
Counting time, monitor      & $T$, $M_j$   & \code{normalization} & yes & yes \\
Integration in radian       & --           & integration path     & yes & n/a \\
Out-of-plane acceptance     & $\dgam$      & \code{acceptance}    & yes & \textbf{no} \\
Lorentz                     & $L_\varphi$ / $L_s$ & \code{geometry} & yes & yes \\
Rod interception            & $\Crod$      & \code{geometry}      & yes & \textbf{no} \\
Polarization                & $P$          & \code{detector}      & yes & yes \\
\quad arm correction        & --           & \code{detector}      & yes & yes \\
Solid angle                 & $\dOm$       & \code{detector}      & \textbf{no} & \textbf{no} \\
Active area / footprint     & $A$          & \code{activearea}    & yes & yes \\
Area correction             & $\Carea$     & \code{geometry}      & slits only & slits only \\
In-plane acceptance         & $\Cdet$      & ---                  & assumed 1 & assumed 1 \\
\bottomrule
\end{tabular}
\end{center}

\noindent
The three entries that distinguish the modes are $\dgam$, $\Crod$ and the
choice of Lorentz factor. Everything else is common, and the two factors
marked \textbf{no} for both modes are the ones whose absence is a deliberate
result rather than an omission.

\section*{References}
\addcontentsline{toc}{section}{References}

\begin{enumerate}[leftmargin=1.6em]
\item E.~Vlieg, \emph{Integrated intensities using a six-circle surface X-ray
  diffractometer}, J.~Appl.~Cryst. \textbf{30} (1997) 532--543.
\item E.~Vlieg, \emph{ANA --- program for the analysis of surface X-ray
  diffraction data}, Appendix~A (z-axis correction-factor table).
\item J.~Drnec, T.~Zhou, S.~Pintea, W.~Onderwaater, E.~Vlieg, G.~Renaud and
  R.~Felici, \emph{Integration techniques for surface X-ray diffraction data
  obtained with a two-dimensional detector}, J.~Appl.~Cryst. \textbf{47}
  (2014) 365--377.
\end{enumerate}

\end{document}
